\input{../tex-inputs/latex-leseansicht-vorspann}

\section[Theodor von Sosnosky an Arthur Schnitzler, 29. 1. 1900]{L04226 Theodor von Sosnosky an Arthur Schnitzler, 29. 1. 1900}
\nopagebreak\mylabel{L04226v}
\rehead{ }\normalsize\beginnumbering\briefempfaengerindex{Schnitzler, Arthur@\textsc{Schnitzler, Arthur}!zzzSosnosky, Theodor von@\emph{von Theodor von Sosnosky}!1900-01-291@{29. 1. 1900}|(be}
\toendnotes[C]{\smallbreak\pagebreak[2]}
\correspDesc{Versand  durch Theodor von Sosnosky am 29. 1. 1900 in Kremsmünster
\newline{}Erhalt  durch Arthur Schnitzler im Zeitraum [30. 1. 1900
                  – 3. 2. 1900?] in Wien}\toendnotes[C]{\smallbreak}
\Standort{DLA, A:Schnitzler, HS.1985.1.4640.}
\physDesc{Brief, 1 Blatt, 3 Seiten, 1930 Zeichen
\newline{}Handschrift: blaue Tinte, deutsche Kurrent
\newline{}Schnitzler: mit rotem Buntstift eine Unterstreichung }\toendnotes[C]{\smallbreak}
\pstart{}{\pb}Sehr geehrter Herr Doktor!\pend\vspace{0.5em}
\pstart
           Empfangen Sie meinen verbindlichſten Dank für Ihre freundliche Einſendung ſowie für
               deren angenehme Form, die mir die Mühe des Kopierens erſpart. Ich erlaube mir, von
               den 5 Gedichten\pwindex{Schnitzler, Arthur 15.\,5.\,1862 Wien – 21.\,10.\,1931 ebd.@\textsc{Schnitzler, Arthur} (15.\,5.\,1862 Wien – 21.\,10.\,1931 ebd.), \emph{Schriftsteller, Mediziner}!Wie wir so still@\strich\emph{Wie wir so still}|pw}\pwindex{Schnitzler, Arthur 15.\,5.\,1862 Wien – 21.\,10.\,1931 ebd.@\textsc{Schnitzler, Arthur} (15.\,5.\,1862 Wien – 21.\,10.\,1931 ebd.), \emph{Schriftsteller, Mediziner}!Anfang vom Ende@\strich\emph{Anfang vom Ende}|pw}\pwindex{Schnitzler, Arthur 15.\,5.\,1862 Wien – 21.\,10.\,1931 ebd.@\textsc{Schnitzler, Arthur} (15.\,5.\,1862 Wien – 21.\,10.\,1931 ebd.), \emph{Schriftsteller, Mediziner}!Ohnmacht@\strich\emph{Ohnmacht}|pw}\pwindex{Schnitzler, Arthur 15.\,5.\,1862 Wien – 21.\,10.\,1931 ebd.@\textsc{Schnitzler, Arthur} (15.\,5.\,1862 Wien – 21.\,10.\,1931 ebd.), \emph{Schriftsteller, Mediziner}!Morgenandacht@\strich\emph{Morgenandacht}|pw}\pwindex{deutsche Lyrik des 19. Jahrhunderts. Eine poetische Revue@\emph{Die deutsche Lyrik des 19. Jahrhunderts. Eine poetische Revue}|pw} 3 zu behalten: Das namenloſe\pwindex{Schnitzler, Arthur 15.\,5.\,1862 Wien – 21.\,10.\,1931 ebd.@\textsc{Schnitzler, Arthur} (15.\,5.\,1862 Wien – 21.\,10.\,1931 ebd.), \emph{Schriftsteller, Mediziner}!Wie wir so still@\strich\emph{Wie wir so still}|pwv}, »Anfang vom
                  Ende\pwindex{Schnitzler, Arthur 15.\,5.\,1862 Wien – 21.\,10.\,1931 ebd.@\textsc{Schnitzler, Arthur} (15.\,5.\,1862 Wien – 21.\,10.\,1931 ebd.), \emph{Schriftsteller, Mediziner}!Anfang vom Ende@\strich\emph{Anfang vom Ende}|pw}« und »Ohnmacht\pwindex{Schnitzler, Arthur 15.\,5.\,1862 Wien – 21.\,10.\,1931 ebd.@\textsc{Schnitzler, Arthur} (15.\,5.\,1862 Wien – 21.\,10.\,1931 ebd.), \emph{Schriftsteller, Mediziner}!Ohnmacht@\strich\emph{Ohnmacht}|pw}«. Bezüglich des letztgenannten\pwindex{Schnitzler, Arthur 15.\,5.\,1862 Wien – 21.\,10.\,1931 ebd.@\textsc{Schnitzler, Arthur} (15.\,5.\,1862 Wien – 21.\,10.\,1931 ebd.), \emph{Schriftsteller, Mediziner}!Ohnmacht@\strich\emph{Ohnmacht}|pwv} bin ich zwar
               ganz Ihrer Anſicht, daß »Morgenandacht\pwindex{Schnitzler, Arthur 15.\,5.\,1862 Wien – 21.\,10.\,1931 ebd.@\textsc{Schnitzler, Arthur} (15.\,5.\,1862 Wien – 21.\,10.\,1931 ebd.), \emph{Schriftsteller, Mediziner}!Morgenandacht@\strich\emph{Morgenandacht}|pw}«
               beſſer ſei; dennoch ziehe ich vor, jenes\pwindex{Schnitzler, Arthur 15.\,5.\,1862 Wien – 21.\,10.\,1931 ebd.@\textsc{Schnitzler, Arthur} (15.\,5.\,1862 Wien – 21.\,10.\,1931 ebd.), \emph{Schriftsteller, Mediziner}!Ohnmacht@\strich\emph{Ohnmacht}|pwv} zu behalten, 1.) weil \introOben{}»M.\pwindex{Schnitzler, Arthur 15.\,5.\,1862 Wien – 21.\,10.\,1931 ebd.@\textsc{Schnitzler, Arthur} (15.\,5.\,1862 Wien – 21.\,10.\,1931 ebd.), \emph{Schriftsteller, Mediziner}!Morgenandacht@\strich\emph{Morgenandacht}|pwv}«\introOben{}\strikeout{es} ſehr lange iſt und ich alle Urſache habe, Raum
               zu ſparen, ſind doch ſchon an 100 Dichter mit gegen 400 Gedichten in dem Buche\pwindex{deutsche Lyrik des 19. Jahrhunderts. Eine poetische Revue@\emph{Die deutsche Lyrik des 19. Jahrhunderts. Eine poetische Revue}|pwv}, darunter leider sehr
               lange, die ich nicht umgehen konnte wie \label{K_L04226-1v}\edtext{»\textsc{Heideschenke\pwindex{Lenau, Nikolaus 13.\,8.\,1802 Lenauheim – 22.\,8.\,1850 Wien@\textsc{Lenau, Nikolaus} (13.\,8.\,1802 Lenauheim – 22.\,8.\,1850 Wien), \emph{Schriftsteller, Librettist}!Heideschenke@\strich\emph{Die Heideschenke}|pw}}«, {[}»{]}Blumen Rache\pwindex{Freiligrath, Hermann Ferdinand 1810 Detmold – 1876@\textsc{Freiligrath, Hermann Ferdinand} (1810 Detmold – 1876), \emph{Journalist, Übersetzer, Revolutionär}!Blumen Rache@\strich\emph{Der Blumen Rache}|pw}«}{\lemma{\textnormal{\emph{»Heideschenke«, … Rache«}}}\Cendnote{\textnormal{Die Ballade \emph{Die
                     Heideschenke}\pwindex{Lenau, Nikolaus 13.\,8.\,1802 Lenauheim – 22.\,8.\,1850 Wien@\textsc{Lenau, Nikolaus} (13.\,8.\,1802 Lenauheim – 22.\,8.\,1850 Wien), \emph{Schriftsteller, Librettist}!Heideschenke@\strich\emph{Die Heideschenke}|pwk} von Nikolaus Lenau\pwindex{Lenau, Nikolaus 13.\,8.\,1802 Lenauheim – 22.\,8.\,1850 Wien@\textsc{Lenau, Nikolaus} (13.\,8.\,1802 Lenauheim – 22.\,8.\,1850 Wien), \emph{Schriftsteller, Librettist}|pwk} hat
                  35, \emph{Der Blumen Rache}\pwindex{Freiligrath, Hermann Ferdinand 1810 Detmold – 1876@\textsc{Freiligrath, Hermann Ferdinand} (1810 Detmold – 1876), \emph{Journalist, Übersetzer, Revolutionär}!Blumen Rache@\strich\emph{Der Blumen Rache}|pwk} von Ferdinand Freiligrath\pwindex{Freiligrath, Hermann Ferdinand 1810 Detmold – 1876@\textsc{Freiligrath, Hermann Ferdinand} (1810 Detmold – 1876), \emph{Journalist, Übersetzer, Revolutionär}|pwk} 20 Strophen, \emph{Die deutsche Lyrik des 19. Jahrhunderts. Eine
                        poetische Revue}\pwindex{deutsche Lyrik des 19. Jahrhunderts. Eine poetische Revue@\emph{Die deutsche Lyrik des 19. Jahrhunderts. Eine poetische Revue}|pwk}. Hg. Theodor von
                        Sosnosky\pwindex{Sosnosky, Theodor von 4.\,1.\,1866 Budapest – 3.\,2.\,1943 Wien@\textsc{Sosnosky, Theodor von} (4.\,1.\,1866 Budapest – 3.\,2.\,1943 Wien), \emph{Schriftsteller}|pwk}, Stuttgart: \emph{J. G. Cotta’sche
                        Buchhandlung Nachfolger}{ }1901, S. 97–103 und 152–155.}}}\label{K_L04226-1}. {\pb}2)
               iſt die erotische Situation deutlicher als \label{K_L04226-2v}\edtext{\textsc{Cotta\orgindex{J.G. Cotta’sche Buchhandlung Nachfolger@J.G. Cotta’sche Buchhandlung Nachfolger|pw}}}{\lemma{\textnormal{\emph{Cotta}}}\Cendnote{\textnormal{Die \emph{Cotta’sche Buchhandlung}\orgindex{J.G. Cotta’sche Buchhandlung Nachfolger@J.G. Cotta’sche Buchhandlung Nachfolger|pwk} wurde nach dem Tod Carl von Cottas\pwindex{Cotta, Carl von 6.\,1.\,1835 – 15.\,9.\,1888@\textsc{Cotta, Carl von} (6.\,1.\,1835 – 15.\,9.\,1888)|pwk} von der Verlegerfamilie um Adolf Kröner\pwindex{Kröner, Adolf 26.\,5.\,1836 Stuttgart – 29.\,1.\,1911@\textsc{Kröner, Adolf} (26.\,5.\,1836 Stuttgart – 29.\,1.\,1911), \emph{Verleger, Musiker, Buchhändler}|pwk} besessen und geleitet. Im Jahr
                     1900 führte Robert Kröner\pwindex{Kröner, Robert 10.\,10.\,1869 Stuttgart – 6.\,1.\,1945 Kirchheim unter Teck@\textsc{Kröner, Robert} (10.\,10.\,1869 Stuttgart – 6.\,1.\,1945 Kirchheim unter Teck), \emph{Verleger, Geschäftsführer}|pwk}
                  die Geschäfte.}}}\label{K_L04226-2} lieb sein könnte, ſo daß ich früchten müße, das Gedichte\pwindex{Schnitzler, Arthur 15.\,5.\,1862 Wien – 21.\,10.\,1931 ebd.@\textsc{Schnitzler, Arthur} (15.\,5.\,1862 Wien – 21.\,10.\,1931 ebd.), \emph{Schriftsteller, Mediziner}!Morgenandacht@\strich\emph{Morgenandacht}|pwv} würde mir geſtrichen
               werden.\pend
           
\pstart
           Mehr als 3 Gedichte konnte ich nicht leicht nehmen, da ich von Autoren, die noch kein
               lyrisches \uline{Buch} veröffentlich haben, in der Regel nur
               ein, zwei Gedichte \strikeout{d} genommen habe, bei Ihnen, Herr
               Doktor, als einer litterarisch besonders intereſſanten Persönlichkeit, ohnehin eine
               Ausnahme mache, auch inſofern, als Ihre Gedichte\pwindex{Schnitzler, Arthur 15.\,5.\,1862 Wien – 21.\,10.\,1931 ebd.@\textsc{Schnitzler, Arthur} (15.\,5.\,1862 Wien – 21.\,10.\,1931 ebd.), \emph{Schriftsteller, Mediziner}!Wie wir so still@\strich\emph{Wie wir so still}|pw}\pwindex{Schnitzler, Arthur 15.\,5.\,1862 Wien – 21.\,10.\,1931 ebd.@\textsc{Schnitzler, Arthur} (15.\,5.\,1862 Wien – 21.\,10.\,1931 ebd.), \emph{Schriftsteller, Mediziner}!Anfang vom Ende@\strich\emph{Anfang vom Ende}|pw}\pwindex{Schnitzler, Arthur 15.\,5.\,1862 Wien – 21.\,10.\,1931 ebd.@\textsc{Schnitzler, Arthur} (15.\,5.\,1862 Wien – 21.\,10.\,1931 ebd.), \emph{Schriftsteller, Mediziner}!Ohnmacht@\strich\emph{Ohnmacht}|pw} in der ganzen Sa{\geminationm}lung\pwindex{deutsche Lyrik des 19. Jahrhunderts. Eine poetische Revue@\emph{Die deutsche Lyrik des 19. Jahrhunderts. Eine poetische Revue}|pwv} die
               einzigen ſein werden, die nicht schon vorher gedruckt erſchienen ſind.\pend
           
\pstart
           Mir perſönlich gefällt »\textsc{Anfang v Ende\pwindex{Schnitzler, Arthur 15.\,5.\,1862 Wien – 21.\,10.\,1931 ebd.@\textsc{Schnitzler, Arthur} (15.\,5.\,1862 Wien – 21.\,10.\,1931 ebd.), \emph{Schriftsteller, Mediziner}!Anfang vom Ende@\strich\emph{Anfang vom Ende}|pw}}« weitaus am besten, an und für{ }ſich ſchon und weil es in seiner »\textsc{Anatol\pwindex{Schnitzler, Arthur 15.\,5.\,1862 Wien – 21.\,10.\,1931 ebd.@\textsc{Schnitzler, Arthur} (15.\,5.\,1862 Wien – 21.\,10.\,1931 ebd.), \emph{Schriftsteller, Mediziner}!Anatol@\strich\emph{Anatol}|pw}}-Sti{\geminationm}ung« für Ihre Eigenart das
               charakteriſtiſcheſte iſt.\pend
           
\pstart
           Wollen Sie das \label{K_L04226-3v}\edtext{namenlose Gedicht\pwindex{Schnitzler, Arthur 15.\,5.\,1862 Wien – 21.\,10.\,1931 ebd.@\textsc{Schnitzler, Arthur} (15.\,5.\,1862 Wien – 21.\,10.\,1931 ebd.), \emph{Schriftsteller, Mediziner}!Wie wir so still@\strich\emph{Wie wir so still}|pw}}{\lemma{\textnormal{\emph{namenlose Gedicht}}}\Cendnote{\textnormal{Schnitzler kam dem Vorschlag nicht nach: Das
                     Gedicht\pwindex{Schnitzler, Arthur 15.\,5.\,1862 Wien – 21.\,10.\,1931 ebd.@\textsc{Schnitzler, Arthur} (15.\,5.\,1862 Wien – 21.\,10.\,1931 ebd.), \emph{Schriftsteller, Mediziner}!Wie wir so still@\strich\emph{Wie wir so still}|pwkv} trägt als
                  Überschrift den Anfang des ersten Verses: »Wie wir so still {\dots}«.}}}\label{K_L04226-3} nicht taufen? Die erste Zeile als Überschrift {\pb}ſcheint mir nicht
               recht geeignet; vielleicht: »\textsc{Stilles Einverständniss}« oder
               »Die Ahnungslosen« oder „\textsc{Wenn Sie wüssten...!}«\pend
           
\pstart
           Die 2 \label{K_L04226-4v}\edtext{nicht verwendeten Gedichte\pwindex{Schnitzler, Arthur 15.\,5.\,1862 Wien – 21.\,10.\,1931 ebd.@\textsc{Schnitzler, Arthur} (15.\,5.\,1862 Wien – 21.\,10.\,1931 ebd.), \emph{Schriftsteller, Mediziner}!Morgenandacht@\strich\emph{Morgenandacht}|pw}\pwindex{deutsche Lyrik des 19. Jahrhunderts. Eine poetische Revue@\emph{Die deutsche Lyrik des 19. Jahrhunderts. Eine poetische Revue}|pw}}{\lemma{\textnormal{\emph{nicht … Gedichte}}}\Cendnote{\textnormal{Welches Gedicht\pwindex{Schnitzler, Arthur 15.\,5.\,1862 Wien – 21.\,10.\,1931 ebd.@\textsc{Schnitzler, Arthur} (15.\,5.\,1862 Wien – 21.\,10.\,1931 ebd.), \emph{Schriftsteller, Mediziner}!?? [Gedicht von Arthur Schnitzler, das Theodor von Sosnosky nicht in seine Lyrikanthologie aufnimmt]@\strich\emph{?? [Gedicht von Arthur Schnitzler, das Theodor von Sosnosky nicht in seine Lyrikanthologie aufnimmt]}|pwkv} neben \emph{Morgenandacht}\pwindex{Schnitzler, Arthur 15.\,5.\,1862 Wien – 21.\,10.\,1931 ebd.@\textsc{Schnitzler, Arthur} (15.\,5.\,1862 Wien – 21.\,10.\,1931 ebd.), \emph{Schriftsteller, Mediziner}!Morgenandacht@\strich\emph{Morgenandacht}|pwk} zurück an Schnitzler
                  ging, ist nicht bekannt.}}}\label{K_L04226-4}{ }ſend ich anbei zurück. Die andern würde ich in dem
               Falle, als \strikeout{ich} meine Anthologie\pwindex{deutsche Lyrik des 19. Jahrhunderts. Eine poetische Revue@\emph{Die deutsche Lyrik des 19. Jahrhunderts. Eine poetische Revue}|pwv} wider Erwarten \textsc{Cottas\orgindex{J.G. Cotta’sche Buchhandlung Nachfolger@J.G. Cotta’sche Buchhandlung Nachfolger|pw}} Sanktion nicht erhalten sollte, natürlich auch zurückſenden, hoffe aber, daß es
               nicht hiezu ko{\geminationm}en wird.\pend
           
\pstart
           Nochmals für Ihr freundliches Entgegenkommen beſtens dankend zeichne ich, Herr
               Doktor, mit vorzügliche{[}r{]}{\\[\baselineskip]}Hochachtung \spacefill\mbox{Theodor vSosnosky}\pend
           \leftskip=0em{}
\pstart
           \noindent{}zt. \textsc{Kremsmünster\oindex{Kremsmünster@\textbf{Kremsmünster}, \emph{Verwaltungsgebiet}|pw}}\pend
           
\pstart
           29./I 900\pend
           \selectlanguage{ngerman}\endnumbering\briefempfaengerindex{Schnitzler, Arthur@\textsc{Schnitzler, Arthur}!zzzSosnosky, Theodor von@\emph{von Theodor von Sosnosky}!1900-01-291@{29. 1. 1900}|)be}\mylabel{L04226h}
\begin{anhang}
\end{anhang}\newcommand{\dateiname}{L04226}\newcommand{\titel}{Theodor von Sosnosky an Arthur Schnitzler, 29. 1. 1900}\newcommand{\editorInnen}{Herausgegeben von Jahnke, SelmaMüller, Martin Anton}\input{../tex-inputs/latex-leseansicht-abspann}
